% !TEX root = ../Hamiltonian_of_mean_force.tex
\section{Dressing invariants for \texorpdfstring{$N$}{N}-level systems\label{sec:general_N}}

\subsection{Block decomposition and first-order invariants\label{subsec:block_decomp}}

For a general $N$-level system with $H_Q\ket{i}=E_i\ket{i}$, the coupling $f$ is an $N\times N$
Hermitian matrix and the influence exponent $\Delta$ is determined by
Eqs.~\eqref{eq:Delta_matrix_elements}--\eqref{eq:Delta_selfenergy}. To extract the physical
invariants, it is instructive to first restrict attention to a single pair of levels
$(i,j)$ with $f_{ij}\neq 0$.

The $(i,j)$ block of $f$ defines a $2\times 2$ sub-problem with energy gap $\Delta_{ij}=E_i-E_j$.
Within this block, the coupling has the same algebraic structure as the qubit case: the traceless
part of the restricted influence exponent satisfies $(M_{ij})^2=\chi_{ij}^2\,\mathbb{I}_{2\times 2}$
with
\begin{equation}
    \chi_{ij}^2 = \lvert f_{ij}\rvert^2\,R(\Delta_{ij}).
    \label{eq:chi_ij_first_order}
\end{equation}
Each non-zero off-diagonal coupling $f_{ij}$ thus generates an independent \emph{first-order}
dressing invariant $\chi_{ij}$, determined by the single bath object $R(\Delta_{ij})$ evaluated
at the corresponding energy gap. At this level the problem decomposes as a direct sum over pairs:
\begin{equation}
    \Delta_{ii}=\sum_j\lvert f_{ij}\rvert^2 R(\Delta_{ij}),
    \label{eq:Delta_diag_sum}
\end{equation}
and the $N$ diagonal elements of $\Delta$ are fixed by the $N(N-1)/2$ first-order invariants
$\{\chi_{ij}\}$. Interactions between different pairs enter only through the off-diagonal
elements $\Delta_{i\ell}$ with $i\neq\ell$, which are the subject of the next subsection.

\subsection{Recursive structure via Cayley--Hamilton\label{subsec:cayley_hamilton}}

The off-diagonal elements of $\Delta$ involve open chains $i\to j\to\ell$ through the intermediate
coupling vertex $j$:
\begin{equation}
    \Delta_{i\ell}
    = \sum_j f_{ij}\,f_{j\ell}\,G^{>}(\Delta_{ij},\Delta_{j\ell}).
    \label{eq:Delta_iell_chain}
\end{equation}
These are \emph{second-order} contributions: each term is a product of two coupling matrix elements
and two bath propagators $G^{>}$, connecting three levels. Continuing the expansion, the matrix
$\Delta^2$ (and its off-diagonal elements) involves chains $i\to j\to k\to\ell$ through two
intermediate vertices, giving \emph{third-order} invariants built from three $G^{>}$ factors, and so on.

The Cayley--Hamilton theorem provides the exact bound on this recursion. For any $N\times N$ matrix
$\Delta$, the characteristic polynomial
\begin{equation}
    p(\lambda)\equiv\det(\lambda\,\mathbb{I}-\Delta)
    =\lambda^N - c_{N-1}\lambda^{N-1}-\cdots-c_1\lambda-c_0
    \label{eq:char_poly}
\end{equation}
is satisfied by $\Delta$ itself: $p(\Delta)=0$. Consequently, any polynomial in $\Delta$ of degree
$\geq N$ can be reduced using $p$, and the full operator algebra generated by $\Delta$ closes at
polynomial degree $N-1$. The \emph{complete set of independent dressing invariants} is therefore
the $N$ coefficients $\{c_k\}_{k=0}^{N-1}$ of the characteristic polynomial---equivalently, the
elementary symmetric polynomials in the eigenvalues $\{\lambda_i\}$ of $\Delta$:
\begin{equation}
    c_k = (-1)^{N-k}\sum_{i_1<\cdots<i_{N-k}}\lambda_{i_1}\cdots\lambda_{i_{N-k}}.
    \label{eq:char_coeffs}
\end{equation}
(One of these, $c_0=\det\Delta$, fixes the normalisation and is determined by the others up to
a sign.) The mean-force Hamiltonian $H_{\mathrm{MF}}=-\beta^{-1}\log(e^{-\beta H_Q}e^{\Delta})$
is therefore a function of $H_Q$ and the $N-1$ independent scalar invariants $\{c_k\}$ alone.

The qubit is the simplest instance of this structure: $N=2$, the characteristic polynomial of
$\Delta$ has one non-trivial coefficient $c_1=\mathrm{tr}\,\Delta=2\Delta_0$, and the traceless
invariant $\chi^2=\Delta_z^2+\Sigma_+\Sigma_-$ is precisely the remaining independent combination
$c_0=\det(M)$ of the traceless part $M$. The single invariant $\chi$ exhausts the full algebraic
content of the qubit $\Delta$, which is why the exact answer collapses to a one-parameter family.

\subsection{Symmetry constraints and rank collapse\label{subsec:symmetry_rank}}

In practice, the coupling operator $f$ is far from a generic $N\times N$ Hermitian matrix: it is
constrained by the physical symmetries of the problem (lattice geometry, particle number, parity,
etc.). These constraints reduce the rank $r$ of $f$, which in turn controls how many of the
$N-1$ Cayley--Hamilton invariants are genuinely independent.

If $f$ has rank $r\leq N$, then $f=\sum_{a=1}^r\ket{v_a}\bra{v_a}$ (schematically), and the
influence exponent is a sum of $r^2$ bilinear terms $f_{ij}f_{jk}$. The matrix $\Delta$ has rank
at most $r^2$ in the eigenbasis of $H_Q$, and its characteristic polynomial has at most $r$
non-trivial roots, giving at most $r$ independent dressing invariants. In the extreme rank-1 limit
$f=\ket{v}\bra{v}$, every matrix element $f_{ij}f_{jk}$ factorises as $v_i\bar v_j\cdot\bar v_j v_k$,
and $\Delta$ collapses to a single-channel object $\Delta\approx c\,f^2$ with one invariant $\chi$,
regardless of $N$. The qubit is therefore not special: it is the generic representative of the
rank-1 universality class.

More broadly, the number of independent order parameters is
\begin{equation}
    \#\{\text{invariants}\} = \mathrm{rank}(f)\;\leq\;N-1,
    \label{eq:n_invariants}
\end{equation}
with the bound saturated only when $f$ couples all $N$ levels democratically. Physical models
typically sit far below this bound: the transverse-field Ising model has $r=1$ (one invariant
per site), the spin-$s$ system has $r=2s$, and so on.

The structure of the invariants is also highly constrained by symmetry. If $f$ commutes with a
symmetry group $G$ of $H_Q$, then $\Delta$ inherits the same symmetry, and the characteristic
polynomial of $\Delta$ decomposes into irreducible blocks of $G$: one independent invariant per
irreducible sector. The Cayley--Hamilton recursion therefore respects the symmetry decomposition
at every order, and the order parameters organise into representations of $G$.

For practical computation at general $N$: (i) diagonalise $H_Q$ to find the energy gaps
$\{\Delta_{ij}\}$; (ii) compute the $r(r+1)/2$ independent elements of $G^{>}$ via
Eq.~\eqref{eq:Ggt_closed_form}; (iii) assemble $\Delta$ from Eq.~\eqref{eq:Delta_matrix_elements};
(iv) evaluate $H_{\mathrm{MF}}$ via the matrix logarithm $-\beta^{-1}\log(e^{-\beta H_Q}e^{\Delta})$.
No further approximation is required; the Cayley--Hamilton theorem guarantees the result is an
exact polynomial in $\Delta$ (and hence in $f^2$) of degree at most $N-1$, with scalar
coefficients drawn entirely from the bath spectral functions $\mathcal{K}(\omega)$ and $R(\omega)$.
